We compare four learning-rate treatments of one hand-written AdamW on a \rhval{const/n_layers}-layer, $d=\rhval{const/dmodel}$ character-level GPT (\rhval{const/n_params} parameters) trained for \rhval{const/steps} steps on Tiny Shakespeare: constant, linear warmup plus cosine decay, step decay, and a reimplemented Schedule-Free AdamW. Each is run at five peak learning rates ($3\times10^{-4}$ to $3\times10^{-2}$) with \rhval{const/n_seeds} seeds, and we define peak-LR sensitivity as the worst-minus-best validation loss over the grid. At the best of the three main peak rates, warmup+cosine (\rhval{main/warmup+cosine/lr3e-3/val_loss/mean:2} nats at $3\times10^{-3}$) beats constant LR (\rhval{main/constant/lr1e-3/val_loss/mean:2} at $10^{-3}$; Welch $p=\rhval{cmp/main/constant/sens/best3/welch_p}$), but Schedule-Free AdamW is better still (\rhval{main/schedule-free-adamw/lr1e-2/val_loss/mean:2} at $10^{-2}$). Our hypothesis that Schedule-Free is less sensitive to the peak LR than warmup+cosine is refuted: over the three main rates its sensitivity is \rhval{main/schedule-free-adamw/sens/sens3/mean:2} nats against \rhval{main/warmup+cosine/sens/sens3/mean:2} for cosine. Most of the large gaps between schedules at high peak rates are a warmup effect rather than a schedule effect: constant and step decay as specified have no warmup and reach \rhval{main/constant/lr1e-2/val_loss/mean:2} and \rhval{main/step-decay/lr1e-2/val_loss/mean:2} at $10^{-2}$, while giving constant LR a 100-step warmup removes most of the gap (\rhval{abl_warmup/constant-warmup-100/lr1e-2/val_loss/mean:2}). With three seeds, one model size and one budget, these are small-scale findings, and the schedule-free result depends on our own implementation and untuned warmup.
